% =====================================================================
%  equations_review.tex  --  SCRATCH / REVIEW COPY (not compiled into the thesis)
% ---------------------------------------------------------------------
%  Purpose: a single place to hand-edit every equation in the thesis.
%  Workflow:
%    1. Edit any equation below to taste (wording, symbols, layout).
%    2. Tell Claude "propagate equations_review.tex" and it will read this
%       file and copy each block back to the SOURCE location noted in the
%       % SOURCE: comment above it.
%  Conventions enforced here (see MATH_STYLE.md):
%    - multi-letter operators  -> \operatorname{...}
%    - loss / metric subscripts -> \text{...}
%    - transpose -> ^\top ; indicator -> \mathbf{1}[...] ; aux weight -> \lambda_{\text{aux}}
%  You can compile this file on its own to preview (pdflatex equations_review.tex).
% =====================================================================
\documentclass[11pt]{article}
\usepackage{amsmath,amssymb,amsbsy}
\allowdisplaybreaks
\begin{document}

% =====================================================================
% CHAPTER 2 -- Literature Review
% =====================================================================

% ---------------------------------------------------------------------
% SOURCE: Chapter/2_Literature_review.tex  (LSTM gates, align block ~L97-104)
% READING: forget/input/output gates, candidate cell state, cell update, hidden update.
% NOTE: currently the first two rows carry \notag so only the last row is numbered.
%       If you want one number for the whole block, use the single-equation form below instead.
\begin{align}
f_t &= \sigma(W_f x_t + U_f h_{t-1} + b_f), &
i_t &= \sigma(W_i x_t + U_i h_{t-1} + b_i), \notag\\
o_t &= \sigma(W_o x_t + U_o h_{t-1} + b_o), &
\tilde{c}_t &= \tanh(W_c x_t + U_c h_{t-1} + b_c), \notag\\
c_t &= f_t \odot c_{t-1} + i_t \odot \tilde{c}_t, &
h_t &= o_t \odot \tanh(c_t),
\end{align}
% (alt: single number for the block)
% \begin{equation}
% \begin{aligned}
% f_t &= \sigma(W_f x_t + U_f h_{t-1} + b_f), &
% i_t &= \sigma(W_i x_t + U_i h_{t-1} + b_i), \\
% o_t &= \sigma(W_o x_t + U_o h_{t-1} + b_o), &
% \tilde{c}_t &= \tanh(W_c x_t + U_c h_{t-1} + b_c), \\
% c_t &= f_t \odot c_{t-1} + i_t \odot \tilde{c}_t, &
% h_t &= o_t \odot \tanh(c_t).
% \end{aligned}
% \end{equation}
% where $\sigma$ is the logistic sigmoid and $\odot$ is the element-wise product.

% ---------------------------------------------------------------------
% SOURCE: Chapter/2_Literature_review.tex  (scaled dot-product attention ~L120-122)
% READING: each output position is a softmax-weighted combination of all value rows.
\begin{equation}
\operatorname{Attention}(Q,K,V) = \operatorname{softmax}\!\left(\frac{QK^\top}{\sqrt{d_k}}\right) V,
\end{equation}
% with $d_k$ the key dimension.

% ---------------------------------------------------------------------
% SOURCE: Chapter/2_Literature_review.tex  (inline ~L139, softmax probabilities)
% READING: class probabilities are the softmax of the logits.
$ p_k = \operatorname{softmax}(z)_k = e^{z_k}/\sum_j e^{z_j} $

% ---------------------------------------------------------------------
% SOURCE: Chapter/2_Literature_review.tex  (cross-entropy ~L143-145)
% READING: negative log-probability of the true class.
\begin{equation}
\mathcal{L}_{\text{CE}} = -\log p_{y},
\end{equation}

% ---------------------------------------------------------------------
% SOURCE: Chapter/2_Literature_review.tex  (inline ~L149, weighted CE)
% READING: cross-entropy scaled by the class weight of the true class.
$ \mathcal{L}_{\text{WCE}} = -\,w_{y}\log p_{y} $

% ---------------------------------------------------------------------
% SOURCE: Chapter/2_Literature_review.tex  (ordinal / squared-EMD loss, label eq:emd ~L155-157)
% READING: mean squared distance between predicted CDF P_k and target step CDF T_k.
% with (inline ~L154):  P_k = \sum_{j\le k} p_j   and   T_k = \mathbf{1}[k \ge y].
\begin{equation}
\mathcal{L}_{\text{ord}} = \frac{1}{K}\sum_{k=0}^{K-1}\bigl(P_k - T_k\bigr)^2 ,
\label{eq:emd}
\end{equation}

% ---------------------------------------------------------------------
% SOURCE: Chapter/2_Literature_review.tex  (inline ~L172, accuracy)
% READING: fraction of correct predictions = trace of confusion matrix / N.
$ \operatorname{Acc}=\frac{1}{N}\sum_i C_{ii} $

% ---------------------------------------------------------------------
% SOURCE: Chapter/2_Literature_review.tex  (macro-accuracy ~L176-178)
% READING: mean of per-class recalls (each class weighted equally).
\begin{equation}
\operatorname{MacroAcc} = \frac{1}{K}\sum_{i=0}^{K-1}\frac{C_{ii}}{\sum_j C_{ij}},
\end{equation}

% ---------------------------------------------------------------------
% SOURCE: Chapter/2_Literature_review.tex  (inline ~L182, MAE)
% READING: average ordinal distance |i-j| over the confusion matrix.
$ \operatorname{MAE}=\frac{1}{N}\sum_{ij}C_{ij}\,|i-j| $

% ---------------------------------------------------------------------
% SOURCE: Chapter/2_Literature_review.tex  (inline ~L187-189, Cohen's kappa)
% READING: chance-corrected agreement; weighted form; unweighted uses W_ij=1 off-diagonal.
$ \kappa = (p_o-p_e)/(1-p_e) $
$ \kappa = 1 - \dfrac{\sum_{ij} W_{ij}C_{ij}}{\sum_{ij} W_{ij}E_{ij}} $

% ---------------------------------------------------------------------
% SOURCE: Chapter/2_Literature_review.tex  (QWK penalty, label eq:qwk ~L192-194)
% READING: quadratic penalty matrix used by quadratic-weighted kappa.
\begin{equation}
W_{ij} = \frac{(i-j)^2}{(K-1)^2},
\label{eq:qwk}
\end{equation}

% ---------------------------------------------------------------------
% SOURCE: Chapter/2_Literature_review.tex  (Spearman rho ~L209-211)
% READING: rank correlation from squared rank differences (no-ties form).
\begin{equation}
\rho = 1 - \frac{6\sum_{i=1}^{n} d_i^{2}}{n\,(n^{2}-1)},
\end{equation}

% ---------------------------------------------------------------------
% SOURCE: Chapter/2_Literature_review.tex  (Kendall tau ~L217-219)
% READING: (concordant - discordant) pairs over the number of pairs C(n,2).
\begin{equation}
\tau = \frac{n_c - n_d}{\binom{n}{2}},
\end{equation}

% ---------------------------------------------------------------------
% SOURCE: Chapter/2_Literature_review.tex  (oracle accuracy ~L226-228)
% READING: fraction of clips that at least one of M configurations gets right.
\begin{equation}
\operatorname{Acc}_{\text{oracle}} = \frac{1}{N}\sum_{n=1}^{N}
  \mathbf{1}\!\left[\,\exists\, m\in\{1,\dots,M\}:\ \hat{y}^{(m)}_n = y_n\,\right],
\end{equation}

% =====================================================================
% CHAPTER 3 -- Methodology
% =====================================================================

% ---------------------------------------------------------------------
% SOURCE: Chapter/3_Methodology.tex  (inverse-frequency class weights ~L243-246)
% READING: weight inversely proportional to class count, normalised to mean 1.
%          n_c = #training samples of class c ; C = #non-empty classes.
\begin{equation}
w_c = \frac{1}{n_c}\cdot\frac{C}{\sum_{k} 1/n_k},
\qquad\text{so that}\quad \frac{1}{C}\sum_c w_c = 1 .
\end{equation}

% ---------------------------------------------------------------------
% SOURCE: Chapter/3_Methodology.tex  (hybrid total loss ~L251-254)
% READING: main-head loss + aux weight * mean of per-stream CE aux losses.
%          z = main logits ; a^{(s)} = stream-s aux logits ; S = #streams (5, or 6 with I3D).
%          \lambda_{\text{aux}} = 0.2 (defined in prose just after the equation).
\begin{equation}
\mathcal{L} = \mathcal{L}_{\text{main}}(z,y) \;+\; \lambda_{\text{aux}}\cdot\frac{1}{S}\sum_{s=1}^{S}
              \mathcal{L}_{\text{CE}}\!\bigl(a^{(s)},y\bigr),
\end{equation}

\end{document}
